\section{RQ1.3 Explanatory Synthesis: Activation-Transfer Size and Compute Distribution}
\label{sec:rq13}

\subsection{Research Question}

This subsection addresses the following sub-research question:

\begin{quote}
\textit{How do activation-transfer size and compute distribution explain the latency and boundary-crossing behavior observed for the decomposition choices evaluated in RQ1.1 and RQ1.2?}
\end{quote}

RQ1.3 is not an additional split benchmark. It is an explanatory synthesis stage that interprets the local findings already reported in RQ1.1 and RQ1.2 using two structural properties of the partition: the size of the intermediate activation tensor transferred at the boundary, and the distribution of computation across the two resulting services. To support this interpretation, RQ1.3 draws on the split benchmark measurements from RQ1.1 and RQ1.2 as primary evidence and uses an independently conducted internal timing experiment on the monolithic ResNet-18 model as supporting evidence for the compute-distribution explanation.

The literature framing for activation-transfer size, compute distribution, and
the ResNet-18 stage structure used as the basis for this synthesis is
discussed in \cref{sec:rw-split-computing,sec:rw-granularity}.

\subsection{Analytical Setup}
\label{sec:rq13-setup}

RQ1.3 draws on three bodies of evidence:

\begin{enumerate}
    \item \textbf{Split benchmark measurements from RQ1.1.} These include the mean end-to-end latencies, boundary-crossing intervals, and activation-transfer sizes for the five coarse conditions (monolithic plus four stage-level splits). The relevant data are summarised in \cref{tab:rq11-results} and the overhead analysis in \cref{sec:rq11-results}.

    \item \textbf{Split benchmark measurements from RQ1.2.} These include the mean end-to-end latencies, boundary-crossing intervals, and activation-transfer sizes for the four refined conditions (monolithic plus three refined splits within the carried-forward region). The relevant data are summarised in \cref{tab:rq12-results} and the overhead analysis in \cref{sec:rq12-results}.

    \item \textbf{Internal timing experiment on monolithic ResNet-18.} A separate instrumented forward-pass profiling experiment was conducted on the same monolithic model under a closely matched single-threaded CPU-stabilised environment. This experiment measured per-stage, per-block, and selected internal operation timings for a single-threaded CPU forward pass, averaged over 200 measured iterations preceded by 50 warmup iterations, using 10 rounds and the same randomisation seed. Functional equivalence was verified (maximum absolute difference of 0.0 against the uninstrumented model), and instrumentation overhead was measured at 0.645~ms (0.88\% of total forward-pass time), confirming that the profiling did not meaningfully distort the timing distribution. This experiment is used as supporting evidence to characterise the internal compute distribution of ResNet-18; it is not a split benchmark and does not introduce new split conditions.
\end{enumerate}

No new split benchmark was conducted for RQ1.3. The explanatory analysis is based entirely on measurements already reported in RQ1.1 and RQ1.2, supplemented by the internal timing experiment described above. To keep this supporting evidence auditable without turning RQ1.3 into a separate benchmark stream, \cref{tab:rq13-config,tab:rq13-internal-summary,fig:rq13-stage-compute-distribution,tab:rq13-topops} provide a compact summary of the internal timing configuration and the specific internal timing results used in the explanatory analysis.

\begin{table}[t]
\centering
\caption{RQ1.3 supporting internal timing configuration}
\label{tab:rq13-config}
\begin{tabular}{ll}
\toprule
\textbf{Setting} & \textbf{Value} \\
\midrule
Model                       & ResNet-18 (pretrained) \\
Device                      & CPU \\
Precision                   & FP32 \\
Profiling mode              & Monolithic internal timing (supporting evidence) \\
Timing scope                & Stage-, block-, and selected operation-level timing \\
Rounds                      & 10 \\
Warmup iterations           & 50 \\
Measured iterations         & 200 \\
Validation inputs           & 5 \\
Validation tolerance        & $1.0 \times 10^{-6}$ \\
CPU stabilisation           & Single-threaded, single-core pinned execution \\
Instrumentation overhead    & 0.645~ms (0.88\%) \\
\bottomrule
\end{tabular}
\end{table}

\begin{table}[t]
\centering
\caption{RQ1.3 supporting internal timing summary used in the explanatory analysis}
\label{tab:rq13-internal-summary}
\begin{tabular}{lrr}
\toprule
\textbf{Unit / Metric} & \textbf{Mean (ms)} & \textbf{\% of model} \\
\midrule
Model total            & 73.753 & 100.00\% \\
Stem                   & 12.028 & 16.31\% \\
Layer1                 & 13.071 & 17.72\% \\
Layer2                 & 11.968 & 16.23\% \\
Layer3                 & 13.689 & 18.56\% \\
Layer4                 & 22.639 & 30.70\% \\
Layer3.1 block         & 7.340  & 9.95\% \\
Classification head    & 0.324  & 0.44\% \\
Convolution ops        & ---    & 81.8\% \\
BatchNorm ops          & ---    & 3.3\% \\
\bottomrule
\end{tabular}
\end{table}

\begin{figure}[t]
\centering
\resizebox{0.92\textwidth}{!}{%
\begin{tikzpicture}[
    font=\small,
    axis/.style={
        thick,
        draw=black!70,
        -{Latex[length=2mm]}
    },
    grid/.style={
        draw=black!12,
        thin
    },
    stagebar/.style={
        draw=blue!65!black,
        fill=blue!18,
        rounded corners=2pt,
        thick
    },
    headbar/.style={
        draw=red!70!black,
        fill=red!10,
        rounded corners=2pt,
        thick
    },
    layerfourbar/.style={
        draw=green!45!black,
        fill=green!18,
        rounded corners=2pt,
        thick
    },
    label/.style={
        font=\scriptsize,
        text=black!75,
        align=right
    },
    vallabel/.style={
        font=\scriptsize\bfseries,
        text=black!75
    },
    note/.style={
        font=\scriptsize,
        text=black!65,
        align=center
    }
]

% x scale: x = percentage * 0.20
% 30.70 -> 6.140
% 18.56 -> 3.712
% 17.72 -> 3.544
% 16.31 -> 3.262
% 16.23 -> 3.246
% 0.44  -> 0.088

% Grid and x-axis labels
\foreach \x/\pct in {
    0/0,
    1/5,
    2/10,
    3/15,
    4/20,
    5/25,
    6/30
} {
    \draw[grid] (\x,0) -- (\x,6.2);
    \node[
        font=\scriptsize,
        text=black!65
    ] at (\x,-0.28) {\pct};
}

% Axes
\draw[axis] (0,0) -- (6.6,0);
\draw[axis] (0,0) -- (0,6.45);

\node[
    font=\small,
    text=black!75
] at (3.3,-0.75) {Share of total forward-pass compute (\%)};

\node[
    rotate=90,
    font=\small,
    text=black!75
] at (-1.85,3.15) {ResNet-18 component};

% Bars
\draw[layerfourbar] (0,5.20) rectangle (6.140,5.70);
\draw[stagebar]     (0,4.40) rectangle (3.712,4.90);
\draw[stagebar]     (0,3.60) rectangle (3.544,4.10);
\draw[stagebar]     (0,2.80) rectangle (3.262,3.30);
\draw[stagebar]     (0,2.00) rectangle (3.246,2.50);
\draw[headbar]      (0,1.20) rectangle (0.088,1.70);

% Y labels
\node[label, anchor=east] at (-0.15,5.45) {\texttt{layer4}};
\node[label, anchor=east] at (-0.15,4.65) {\texttt{layer3}};
\node[label, anchor=east] at (-0.15,3.85) {\texttt{layer1}};
\node[label, anchor=east] at (-0.15,3.05) {Stem};
\node[label, anchor=east] at (-0.15,2.25) {\texttt{layer2}};
\node[label, anchor=east] at (-0.15,1.45) {Classification head};

% Value labels
\node[vallabel, anchor=west] at (6.24,5.45) {30.70\%};
\node[vallabel, anchor=west] at (3.82,4.65) {18.56\%};
\node[vallabel, anchor=west] at (3.65,3.85) {17.72\%};
\node[vallabel, anchor=west] at (3.37,3.05) {16.31\%};
\node[vallabel, anchor=west] at (3.35,2.25) {16.23\%};
\node[vallabel, anchor=west] at (0.22,1.45) {0.44\%};

% Annotation for the degenerate downstream segment
\draw[
    red!70!black,
    thick,
    -{Latex[length=1.8mm]}
] (1.55,0.85) -- (0.18,1.20);

\node[
    note,
    text width=6.0cm,
    anchor=north
] at (3.25,0.75) {
    After \texttt{layer4}, only the classification head remains.
    This is below the 10\% compute-degeneracy threshold.
};

\end{tikzpicture}%
}
\caption{Stage-level compute distribution from the RQ1.3 internal timing analysis. The final classification head contributes less than 1\% of total forward-pass compute, while \texttt{layer4} accounts for the largest stage-level share. This explains why \texttt{split\_after\_layer4} is raw-fast but compute-degenerate: almost all meaningful model computation remains before the split.}
\label{fig:rq13-stage-compute-distribution}
\end{figure}

\begin{table}[t]
\centering
\caption{RQ1.3 heaviest internal compute units from the supporting timing analysis}
\label{tab:rq13-topops}
\begin{tabular}{lrr}
\toprule
\textbf{Unit} & \textbf{Mean (ms)} & \textbf{\% of model} \\
\midrule
\texttt{stem.maxpool}   & 7.189 & 9.75\% \\
\texttt{layer4.1.conv1} & 6.039 & 8.19\% \\
\texttt{layer4.0.conv2} & 6.029 & 8.17\% \\
\texttt{layer4.1.conv2} & 5.972 & 8.10\% \\
\texttt{layer3.0.conv2} & 3.605 & 4.89\% \\
\texttt{stem.conv1}     & 3.598 & 4.88\% \\
\texttt{layer3.1.conv1} & 3.576 & 4.85\% \\
\texttt{layer4.0.conv1} & 3.485 & 4.73\% \\
\bottomrule
\end{tabular}
\end{table}

\subsection{Explanatory Findings}
\label{sec:rq13-findings}

The findings are organised around four explanatory observations, each building on the previous one.

\subparagraph{A. Activation-transfer burden explains much of the coarse latency pattern.}

In RQ1.1, the boundary-crossing interval decreased monotonically from the earliest split to the latest: 53.03~ms after \texttt{layer1} (784~KB activation), 39.96~ms after \texttt{layer2} (392~KB), 25.73~ms after \texttt{layer3} (196~KB), and 2.13~ms after \texttt{layer4} (98~KB). This pattern closely tracks the reduction in activation-transfer size as the split point moves deeper into the network. Larger intermediate tensors require more serialisation, transmission, and deserialisation time over the gRPC boundary, and this cost is directly reflected in the boundary-crossing measurements.

This observation is consistent with the broader partitioning literature, which identifies activation size as a primary driver of communication overhead in distributed inference~\cite{neurosurgeon, jointdnn, bottlenet_pp, deeperthings, survey_partitioning}. At the coarse level, the monotonic relationship between activation size and boundary-crossing interval confirms that transfer burden is the dominant explanatory factor for the boundary-crossing cost gradient observed across the four RQ1.1 split points.

\subparagraph{B. Activation size alone does not determine suitability.}

Although \texttt{split\_after\_layer4} achieved the lowest boundary-crossing interval and the lowest raw split latency in RQ1.1 (77.26~ms), it was classified as compute-degenerate because the second service executed only average pooling and the final linear layer --- a residual computation segment of well under 1~ms. This means that minimising activation-transfer size to its smallest value does not, by itself, produce a suitable microservice decomposition. A boundary that pushes nearly all computation to one side leaves the second service structurally empty, which defeats the purpose of decomposition.

The internal timing experiment provides quantitative support for this observation. As shown in \cref{fig:rq13-stage-compute-distribution}, the monolithic forward pass took approximately 73.75~ms on average, of which \texttt{layer4} alone accounted for 30.70\% (approximately 22.64~ms). More importantly, the computation remaining after a \texttt{layer4} split --- average pooling plus the final linear layer --- accounts for only about 0.32~ms in total, or roughly 0.44\% of the full forward pass. This means that a split after \texttt{layer4} places approximately 99.6\% of the total computation on the first service and leaves only about 0.4\% on the second. By contrast, a split after \texttt{layer3} places approximately 68.8\% of compute on the first service and 31.1\% on the second, producing a substantially more balanced distribution.

The supporting timing analysis also shows that compute is concentrated in a small number of heavy internal units rather than being spread uniformly. As shown in \cref{tab:rq13-topops}, the heaviest units are dominated by convolutional operations in later network regions, especially \texttt{layer4}. This makes a very late split structurally problematic: it transfers a small activation tensor, but it also leaves the remote side with almost no substantial computation.

This distinction between raw-fastest and structurally suitable is consistent with literature that treats compute distribution as a first-class partition criterion alongside communication cost~\cite{dnnsplit, dads, pareto_split, bottlefit}.

\subparagraph{C. The RQ1.2 controlled comparison isolates the role of compute placement.}

The most precise evidence for the independent role of compute distribution comes from the RQ1.2 comparison between \texttt{split\_after\_layer3\_block0} and \texttt{split\_after\_layer3}. Both conditions transfer an intermediate tensor of approximately 196~KB. If activation-transfer size were the sole determinant of boundary-crossing cost, these two conditions should exhibit similar boundary-crossing intervals and similar end-to-end latencies. They do not.

\texttt{split\_after\_layer3\_block0} exhibited a boundary-crossing interval of 33.135~ms, while \texttt{split\_after\_layer3} exhibited a boundary-crossing interval of 25.718~ms --- a difference of more than 7.4~ms despite identical activation size. The internal timing data helps explain this gap. The \texttt{layer3.1} block (the second residual block of \texttt{layer3}) accounts for approximately 9.95\% of total monolithic compute. The difference is more directly explained by the amount of computation that remains on the second service. When the split is placed after \texttt{layer3.0}, Service~B executes \texttt{layer3.1} in addition to \texttt{layer4}, pooling, and classification. When the split is placed after the full \texttt{layer3}, that computation remains on Service~A. The internal timing report shows that \texttt{layer3.1} accounts for approximately 7.34~ms of monolithic compute, which closely matches the observed 7.42~ms difference in boundary-crossing interval between the two refined conditions. This provides direct evidence that the boundary-crossing difference is driven primarily by compute placement on the remote side rather than by activation-transfer size, which is identical in both cases.

This comparison also shows why a purely tensor-size-based explanation is incomplete. The two refined conditions share the same activation size, but the internal work performed after the boundary differs materially. In that sense, RQ1.2 isolates compute placement as the main explanatory variable within the carried-forward region.

This observation directly supports the principle established in the partitioning literature that compute distribution is not a secondary consideration but an independent factor in end-to-end split behavior and boundary-crossing cost~\cite{dnnsplit, dads, fine_grained_partitioning}.

\subparagraph{D. Internal timing confirms that compute is unevenly distributed within ResNet-18.}

The internal timing experiment provides a structural explanation for why different boundaries within the same coarse region yield different performance. As visualised in \cref{fig:rq13-stage-compute-distribution}, the per-stage compute shares are approximately: stem 16.31\%, \texttt{layer1} 17.72\%, \texttt{layer2} 16.23\%, \texttt{layer3} 18.56\%, and \texttt{layer4} 30.70\%. Convolution operations account for 81.8\% of total compute, while batch normalisation contributes approximately 3.3\% and operations below 0.1~ms (such as ReLU and residual addition) are individually negligible.

The operation-level support data adds a more concrete view of this distribution. As shown in \cref{tab:rq13-topops}, the heaviest internal units are predominantly convolutional operations, with the largest shares concentrated in \texttt{layer4} and the late \texttt{layer3} region. By contrast, batch normalisation, residual addition, and ReLU operations are individually small and do not serve as meaningful compute anchors by themselves. This confirms that ResNet-18 is not only uneven at the stage level, but also internally dominated by a relatively small number of heavy convolutional units.

Two aspects of this distribution are relevant to the RQ1.1 and RQ1.2 findings. First, \texttt{layer4} is by far the most compute-intensive stage, which explains why a split after \texttt{layer4} is structurally degenerate: it assigns the dominant compute block entirely to the first service. Second, the per-block decomposition within \texttt{layer3} reveals that \texttt{layer3.0} and \texttt{layer3.1} do not contribute equally; the difference in their execution times contributes to the boundary-crossing asymmetry observed between \texttt{split\_after\_layer3\_block0} and \texttt{split\_after\_layer3}. This structural unevenness is consistent with the design of ResNet-18, where each successive stage doubles the channel count and halves the spatial resolution, producing progressively heavier convolution operations in later stages~\cite{resnet}.

It should be noted that the internal timing experiment was conducted on the monolithic model and therefore measures local per-layer execution cost without the serialisation and communication overhead present in the split benchmarks. The timing values are used here to explain the compute-distribution dimension; the actual boundary-crossing cost is a composite of both communication and compute-related factors as measured in the split benchmarks themselves.

The interpretation of these findings and the answer to RQ1.3 are presented
in \cref{sec:analysis-rq13}.